\documentclass[10pt,a4paper]{amsart}
\usepackage[margin=19mm]{geometry}
\usepackage{amsmath,amssymb,mathtools}
\usepackage[scr=rsfs,cal=euler]{mathalfa}
\usepackage{xfrac}
\usepackage[unicode,hidelinks,bookmarks=false]{hyperref}
\newcommand{\ie}{i.e.\ }
\newcommand{\cf}{cf.\ }
\newcommand{\resp}{respectively\ }
\newcommand{\Subsection}{\subsection}
\providecommand{\nobreakdash}{\nobreak\hbox{-}}
\setlength{\emergencystretch}{2em}
\multlinegap0pt
\usepackage{amssymb}
\usepackage[shortlabels]{enumitem}
\usepackage{mathtools}
\usepackage{tikz}
\usepackage{xcolor}
\usepackage{float}
\newtheorem{prop}{Proposition}
\title{The type and cardinality of minimal presentations of \\ numerical semigroups with embedding dimension four}
\author{Kazimierz Chomicz}
\date{Source: arXiv:2609.04000v1 (CC BY 4.0)}
\begin{document}
\maketitle

\noindent\emph{Source and licence.} The type and cardinality of minimal presentations of \\ numerical semigroups with embedding dimension four. Version arXiv:2609.04000v1 (CC BY 4.0), distributed by arXiv under the Creative Commons Attribution 4.0 International licence, http://creativecommons.org/licenses/by/4.0/, as recorded in the arXiv licence metadata for this exact version and shown on the abstract page https://arxiv.org/abs/2609.04000v1. Full licence notice: \texttt{licenses/arxiv-2609.04000-CC-BY-4.0.txt}.\par\smallskip

\noindent\textit{Editorial scope.} Counted unit: the article's prop prop\_type\_suppthm\_odd\_corners (source0.tex lines 1292--1294). The article's complete original proof of it is delivered with it (source0.tex lines 1296--1377, one \texttt{proof} environment of 82 lines). No mathematics is written, repaired or re-derived: the statement and the proof are the article's own lines, in the article's own notation, and no retained line is substituted or re-flowed.  No further source-local context is needed: every cross-reference of every retained line resolves inside the two delivered spans.  One declared adaptation: exactly 6 ancillary strings inside the retained spans are omitted (image-inclusion commands and, where the source repeats a label, the repeated label definitions) and no other line differs from the source; the figure environments, with their placement, captions and labels, are kept, and the package records the omitted strings in fidelity receipts bound to the affected bodies.  No retained line contains a citation, so the wrapper carries no bibliography and no bibliography entry.  The retained lines contain the article's own diagram code, which the wrapper typesets with the standard tikz packages; no image file is delivered and the package declares no asset dependency.  Editorial formatting only, in addition: an AMS \texttt{amsart} preamble that restates the article's own declarations and macros used by the retained lines, namely the environments \texttt{prop} on separate counters, together with the standard packages those lines need.  The retained environments print in this document as Proposition 1 in order of appearance, whereas the article prints the same items with the numbering of its own full document.  The complete native source of the article as distributed in the arXiv e-print for this exact version is bundled unchanged as \texttt{sources/209334/source0.tex}.
\begin{prop}\label{prop_type_suppthm_odd_corners}
    Let $F$ be a cube figure. If there is no $[[a,b,c]] \notin F$ such that $[[a-1,b,c]],[[a,b-1,c]],[[a,b,c-1]] \in F$, then $F$ has at most two odd corners.
\end{prop}

\begin{proof}
    Let $[[a,b,c]]$ be an odd corner of $F$. This implies: $[[a+1,b,0]] \in F$ or $[[a,b+1,0]] \in F$, $[[a+1,0,c]] \in F$ or $[[a,0,c+1]] \in F$, and $[[0,b+1,c]] \in F$ or $[[0,b,c+1]] \in F$. Suppose that $[[a+1,b,0]]$ and $[[a,b+1,0]]$ are both in $F$ --- we claim that this cannot be the case. 

    If $[[a+1,0,c]] \in F$, then there exists $b \geq \lambda > 0$ such that $[[a+1,\lambda-1,c]] \in F$ and $[[a+1,\lambda,c]] \notin F$. Then, there exists $c \geq \mu > 0$ such that $[[a+1,\lambda,\mu-1]] \in F$ and $[[a+1,\lambda,\mu]] \notin F$, as $[[a+1,\lambda,0]] \leq [[a+1,b,0]] \in F$. However, since $[[a+1,\lambda,\mu]] \notin F$ and $[[a,\lambda,\mu]], [[a+1,\lambda-1,\mu]], [[a+1,\lambda,\mu-1]] \in F$, we get a contradiction with the hypothesis. On Figure \ref{fig_type_theorem_1}, we can see an example of such a situation, where $P = (a+1,\lambda,\mu)$. 
    
    If $[[0,b+1,c]] \in F$, we get an analogous contradiction; thus, since $[[a,b,c]]$ is an odd corner and neither $[[a+1,0,c]]$ nor $[[0,b+1,c]]$ is in $F$, we have $[[a,0,c+1]], [[0,b,c+1]] \in F$. Now, there exist $a \geq \lambda > 0$, $b \geq \mu > 0$, such that $[[\lambda,\mu,c+1]] \notin F$ and $[[\lambda-1,\mu,c+1]], [[\lambda,\mu-1,c+1]] \in F$. Since $[[\lambda,\mu,c]] \in F$ also, we get a contradiction with the hypothesis. On Figure \ref{fig_type_theorem_2}, we can see an example of such a situation, where $P = (\lambda,\mu,c+1)$.


\begin{figure}[h]
    \centering
    \begin{minipage}[t]{0.48\textwidth}
        \centering
        
        \caption{}
        \label{fig_type_theorem_1}
    \end{minipage}
    \begin{minipage}[t]{0.48\textwidth}
        \centering
        
        \caption{}
        \label{fig_type_theorem_2}
    \end{minipage}
\end{figure}

    Similarly, we obtain that exactly one of $[[a+1,0,c]]$ and $[[a,0,c+1]]$ is in $F$, and that exactly one of $[[0,b+1,c]]$ and $[[0,b,c+1]]$ is in $F$. 

    Now, we claim that any two distinct odd corners $[[a,b,c]]$ and $[[a',b',c']]$ of $F$ must agree in at least one coordinate. Otherwise, without loss of generality, suppose that $a < a'$, $b < b'$, $c > c'$ (we cannot have $[[a,b,c]] < [[a',b',c']]$, as they are both corners). Then, clearly $[[a+1,b,0]], [[a,b+1,0]] \in F$, since $[[a+1,b,0]], [[a,b+1,0]] < [[a',b',c']] \in F$, which contradicts the claim proved at the beginning.

    
    Thus, without loss of generality, suppose that $a < a'$, $b = b'$, $c > c'$.  Suppose that $[[a'',b'',c'']]$ is a third distinct odd corner of $F$. If $b'' \neq b$, then from the previous paragraph
    we obtain $[[a'',b'',c'']] = [[a,b'',c']]$ or $[[a'',b'',c'']] = [[a',b'',c]]$ --- recall that $a < a'$ and $c > c'$, so we cannot have $a'' = a' = a$ or $c'' = c' = c$. 



    
    Suppose $[[a'',b'',c'']] = [[a',b'',c]]$. If $b'' \geq b' = b$, then $[[a',b'',c]] >[[a',b',c']]$, which contradicts the fact that $[[a',b',c']]$ is a corner of $F$. Hence $ b'' < b' = b$. Now, there exist $a' \geq \lambda > a$, $b \geq \mu > b''$, and $c \geq \nu > c'$, such that $[[\lambda,\mu,\nu]] \notin F$ and $[[\lambda-1,\mu,\nu]],[[\lambda,\mu-1,\nu]],[[\lambda,\mu,\nu-1]] \in F$, which contradicts the hypothesis. On Figure \ref{fig_type_theorem_3}, we can see an example of such a situation, where $P = (\lambda,\mu,\nu)$. 
    
    Suppose $[[a'',b'',c'']] = [[a,b'',c']]$. If $ b'' \leq b'=b $, then $[[a,b'',c']] < [[a,b,c]]$, which contradicts the fact that $ [[a,b'',c']]$ is a corner of $F$. Hence $b'' > b' =b$. Now, since $[[a+1,b,0]] \leq [[a',b',0]] \in F$ and $ [[a,b+1,0]] \leq [[a,b'',0]] = [[a'',b'',0]] \in F$, we have $[[a+1,b,0]], [[a,b+1,0]] \in F$, which contradicts the claim proved at the beginning. On Figure \ref{fig_type_theorem_4}, we can see an example of such a situation. 

\begin{figure}[h]
    \centering
    \begin{minipage}[t]{0.48\textwidth}
        \centering
        
        \caption{}
        \label{fig_type_theorem_3}
    \end{minipage}
    \begin{minipage}[t]{0.48\textwidth}
        \centering
        
        \caption{}
        \label{fig_type_theorem_4}
    \end{minipage}
\end{figure}


    
    Hence $b= b' = b'' $, and without loss of generality, suppose that
    $a < a' < a''$ and $c > c' > c''$. We have $[[a'+1,0,c']] \in F$ or $[[a',0,c'+1]] \in F$, as $[[a',b',c']]$ is an odd corner. 
    
    If $[[a',0,c'+1]] \in F$, then there exist $a' \geq \lambda > a$ and $b \geq \mu > 0$ such that $[[\lambda,\mu,c'+1]] \notin F$ and $[[\lambda-1,\mu,c'+1]], [[\lambda,\mu-1,c'+1]], [[\lambda,\mu,c']] \in F$, which contradicts the hypothesis. On Figure \ref{fig_type_theorem_5}, we can see an example of such a situation, where $P = (\lambda,\mu,c'+1)$.

     If $[[a'+1,0,c']] \in F$, then there exist $a'' \geq \lambda > a'$ and $b \geq \mu > 0$ such that $[[\lambda,\mu,c''+1]] \notin F$ and $[[\lambda-1,\mu,c''+1]], [[\lambda,\mu-1,c''+1]], [[\lambda,\mu,c'']] \in F$, which contradicts the hypothesis. On Figure \ref{fig_type_theorem_6}, we can see an example of such a situation, where $P = (\lambda,\mu,c'+1)$.  \qedhere
     


\begin{figure}[h]
    \centering
    \begin{minipage}[t]{0.48\textwidth}
        \centering
        
        \caption{}
        \label{fig_type_theorem_5}
    \end{minipage}
    \begin{minipage}[t]{0.48\textwidth}
        \centering
        
        \caption{}
        \label{fig_type_theorem_6}
    \end{minipage}
\end{figure}
\end{proof}

\end{document}
